\chapter{Services SOAP}\label{chap:soap}

\section{Conception commune}
Les trois services suivent la même structure, de type \emph{contract-first} par annotations :
\begin{itemize}
  \item une \textbf{interface de service} (SEI) annotée \code{@WebService}, qui définit le contrat et génère le WSDL ;
  \item une \textbf{implémentation} qui applique le contrôle d'accès, la validation et les règles de gestion ;
  \item un \textbf{DAO} JDBC limité aux tables du service ;
  \item des \textbf{DTO JAXB} (\code{@XmlType}) portant les types du schéma XSD ; les dates utilisent le type \code{xs:date}.
\end{itemize}

Chaque opération déclare une faute métier unique, \code{ServiceException}, sérialisée en \code{<detail>} sous la forme
d'un élément \code{ServiceFault} (espace de noms \code{http://stages.universite.com/commun}) :

\begin{table}[htbp]
\centering
\caption{Codes de la faute métier \code{ServiceFault}}\label{tab:faults}
\begin{tabularx}{\textwidth}{@{}lY@{}}
\toprule
\textbf{Code} & \textbf{Signification} \\
\midrule
\code{DONNEES\_INVALIDES} & Échec de validation ; la liste \code{detail} contient \emph{toutes} les erreurs détectées \\
\code{RESSOURCE\_INTROUVABLE} & Matricule ou identifiant inconnu \\
\code{CONFLIT} & Unicité violée (matricule, e-mail, raison sociale) \\
\code{REGLE\_METIER} & Règle de gestion non respectée (suppression interdite, plus de place…) \\
\code{ACCES\_REFUSE} & Rôle ou périmètre de l'appelant insuffisant \\
\code{ERREUR\_TECHNIQUE} & Service tiers indisponible lors d'une composition \\
\bottomrule
\end{tabularx}
\end{table}

Une requête sans en-tête WS-Security valide est rejetée avant d'atteindre le service par une faute standard
\code{wsse:FailedAuthentication} (chapitre~\ref{chap:securite}). Les WSDL et XSD exportés sont fournis dans
\code{docs/wsdl/}.

\section{Service 1 — Gestion des étudiants}
\textbf{Adresse} : \code{/ws/etudiants} \quad \textbf{Espace de noms} : \code{http://stages.universite.com/etudiants}

\begin{table}[htbp]
\centering\small
\caption{Opérations du service Étudiants}
\begin{tabularx}{\textwidth}{@{}lYYl@{}}
\toprule
\textbf{Opération} & \textbf{Entrée} & \textbf{Sortie / règles} & \textbf{Rôles} \\
\midrule
\code{ajouterEtudiant} & \code{Etudiant} & Étudiant créé (statut \code{ACTIF} par défaut) ; matricule et e-mail uniques & Resp. \\
\code{consulterEtudiant} & \code{matricule} & Fiche complète ; un étudiant ne lit que la sienne ; téléphone et e-mail masqués pour une entreprise & Tous \\
\code{modifierEtudiant} & \code{Etudiant} & Fiche mise à jour, mêmes validations ; un étudiant ne modifie que son profil, et seulement e-mail, téléphone, adresse et compétences & Resp., Étud. \\
\code{supprimerEtudiant} & \code{matricule} & \code{boolean} ; interdit si \code{EN\_STAGE} & Resp. \\
\code{listerEtudiants} & \code{filiere}?, \code{statut}? & Liste filtrée triée par nom & Resp., Syst. \\
\code{verifierEtudiant} & \code{matricule} & \code{VerificationEtudiant} : existe, statut, éligible, message & Tous \\
\code{changerStatutEtudiant} & \code{matricule}, \code{statut} & Fiche mise à jour & Resp., Syst. \\
\bottomrule
\end{tabularx}
\end{table}

\textbf{Validation} : matricule au format \code{[A-Z]\{1,3\}[0-9]\{4,10\}} (ex. \code{E2026001}), nom et prénom obligatoires
(80 caractères max), e-mail et téléphone au bon format, niveau dans \{INE1, INE2, INE3\}, statut dans l'énumération.

\section{Service 2 — Gestion des entreprises}
\textbf{Adresse} : \code{/ws/entreprises} \quad \textbf{Espace de noms} : \code{http://stages.universite.com/entreprises}

\begin{table}[htbp]
\centering\small
\caption{Opérations du service Entreprises}
\begin{tabularx}{\textwidth}{@{}lYYl@{}}
\toprule
\textbf{Opération} & \textbf{Entrée} & \textbf{Sortie / règles} & \textbf{Rôles} \\
\midrule
\code{ajouterEntreprise} & \code{Entreprise} & Entreprise créée ; raison sociale unique ; date de référencement automatique & Resp. \\
\code{consulterEntreprise} & \code{entrepriseId} & Fiche complète (secteur, coordonnées, contact) ; une entreprise ne consulte que la sienne, coordonnées masquées pour un étudiant & Tous \\
\code{modifierEntreprise} & \code{Entreprise} & Une entreprise modifie ses coordonnées et son contact, mais ni sa raison sociale, ni son secteur, ni son référencement & Resp., Ent. \\
\code{supprimerEntreprise} & \code{entrepriseId} & Interdit si l'entreprise a des offres actives & Resp. \\
\code{listerEntreprises} & \code{secteurActivite}? & Liste filtrée ; coordonnées des autres entreprises masquées & Tous \\
\code{verifierEntreprise} & \code{entrepriseId} & \code{VerificationEntreprise} : existe, référencée, message & Tous \\
\code{consulterOffresEntreprise} & \code{entrepriseId} & Liste d'\code{Offre} — \emph{composition} : appel SOAP au service Offres & Tous \\
\bottomrule
\end{tabularx}
\end{table}

\section{Service 3 — Gestion des offres de stage}
\textbf{Adresse} : \code{/ws/offres} \quad \textbf{Espace de noms} : \code{http://stages.universite.com/offres}

\begin{table}[htbp]
\centering\small
\caption{Opérations du service Offres}
\begin{tabularx}{\textwidth}{@{}lYYl@{}}
\toprule
\textbf{Opération} & \textbf{Entrée} & \textbf{Sortie / règles} & \textbf{Rôles} \\
\midrule
\code{creerOffre} & \code{Offre} & Statut \code{ACTIVE} ; type de stage obligatoire ; date limite future ; entreprise référencée (appel SOAP Entreprises) & Resp., Ent.\footnotemark \\
\code{consulterOffre} & \code{offreId} & Offre complète & Tous \\
\code{modifierOffre} & \code{Offre} & Places $\geq$ places déjà attribuées ; propriétaire et date de publication figés & Resp., Ent. \\
\code{supprimerOffre} & \code{offreId} & Interdit si des stagiaires ont été retenus & Resp., Ent. \\
\code{listerOffres} & \code{statut}?, \code{entrepriseId}? & Liste triée par date limite ; pour un étudiant, seuls les types ouverts à son niveau & Tous \\
\code{listerOffresDisponibles} & — & Offres actives, non expirées, avec places ; filtrage par niveau pour un étudiant & Tous \\
\code{verifierDisponibilite} & \code{offreId} & \code{DisponibiliteOffre} : active, date dépassée, places, message & Tous \\
\code{verifierCompatibilite} & \code{offreId}, \code{filiere}, \code{niveau}, \code{competences} & \code{Compatibilite} : filière, niveau autorisé par le type de stage, date limite, score de compétences, manques, motifs & Tous \\
\code{reserverPlace} & \code{offreId} & Décrément atomique ; passage à \code{POURVUE} à zéro & Resp., Syst. \\
\code{listerOffresExpirantBientot} & \code{jours} & Offres actives clôturant dans l'intervalle & Tous \\
\bottomrule
\end{tabularx}
\end{table}
\footnotetext{Une entreprise ne peut créer, modifier ou supprimer que ses propres offres.}

\paragraph{Type de stage.} Chaque offre porte un \code{typeStage} obligatoire ; les années d'études autorisées en
découlent (tableau~\ref{tab:types}) et sont renvoyées, à titre informatif, dans \code{niveauxAutorises}. L'entreprise ne
choisit que le type : la règle de l'école est appliquée de façon uniforme par le service. Pour un utilisateur de rôle
étudiant, \code{listerOffres} et \code{listerOffresDisponibles} ne renvoient que les offres dont le type est ouvert à son
niveau, obtenu par composition SOAP auprès du service Étudiants : le filtrage est fait côté serveur.

\begin{table}[htbp]
\centering\small
\caption{Types de stage et niveaux autorisés}\label{tab:types}
\begin{tabularx}{\textwidth}{@{}llY@{}}
\toprule
\textbf{Code} & \textbf{Type de stage} & \textbf{Niveaux autorisés (motif en cas de refus)} \\
\midrule
\code{OBSERVATION} & Stage d'observation (initiation) & INE1 — « Cette offre de stage d'observation est réservée aux étudiants de 1re année (INE1). » \\
\code{OUVRIER\_PFA} & Stage ouvrier / PFA & INE1, INE2 — « … réservée aux étudiants de 1re et 2e année (INE1, INE2). » \\
\code{PFE} & Stage PFE & INE3 — « Cette offre de stage PFE est réservée aux étudiants de 3e année (INE3). » \\
\bottomrule
\end{tabularx}
\end{table}

\paragraph{Compatibilité « intelligente ».} \code{verifierCompatibilite} combine trois critères \emph{bloquants} —
filière ciblée (liste vide = toutes), niveau de l'étudiant parmi ceux qu'ouvre le type de stage, et date limite non
dépassée — et un \textbf{score de
compétences} indicatif : pourcentage des compétences requises que possède l'étudiant, comparaison insensible à la casse
et aux accents. Le score est conservé avec la candidature et présenté à l'entreprise pour l'aider à trier les dossiers.

\paragraph{Réservation sans sur-attribution.} \code{reserverPlace} exécute
\code{UPDATE offre SET places\_disponibles = places\_disponibles - 1 … WHERE id = ? AND places\_disponibles > 0} :
la condition dans la clause \code{WHERE} rend l'opération atomique, même si deux acceptations sont simultanées.

\section{Extrait de WSDL}
\begin{lstlisting}[language=XML, caption={Extrait de \code{etudiants.wsdl} : opération \code{verifierEtudiant}}]
<message name="verifierEtudiant">
  <part name="parameters" element="tns:verifierEtudiant"/>
</message>
<message name="ServiceException">
  <part name="fault" element="ns1:ServiceFault"/>
</message>
<portType name="EtudiantService">
  <operation name="verifierEtudiant">
    <input  message="tns:verifierEtudiant"/>
    <output message="tns:verifierEtudiantResponse"/>
    <fault  message="tns:ServiceException" name="ServiceException"/>
  </operation>
</portType>
<binding name="EtudiantPortBinding" type="tns:EtudiantService">
  <soap:binding transport="http://schemas.xmlsoap.org/soap/http" style="document"/>
  <operation name="verifierEtudiant">
    <soap:operation soapAction=""/>
    <input><soap:body use="literal"/></input>
    <output><soap:body use="literal"/></output>
    <fault name="ServiceException"><soap:fault name="ServiceException" use="literal"/></fault>
  </operation>
</binding>
<service name="EtudiantService">
  <port name="EtudiantPort" binding="tns:EtudiantPortBinding">
    <soap:address location="http://localhost:8090/ws/etudiants"/>
  </port>
</service>
\end{lstlisting}

\begin{lstlisting}[language=XML, caption={Type XSD \code{VerificationEtudiant} généré}]
<xs:complexType name="VerificationEtudiant">
  <xs:sequence>
    <xs:element name="matricule" type="xs:string" minOccurs="0"/>
    <xs:element name="existe"    type="xs:boolean"/>
    <xs:element name="statut"    type="xs:string" minOccurs="0"/>
    <xs:element name="eligible"  type="xs:boolean"/>
    <xs:element name="message"   type="xs:string" minOccurs="0"/>
  </xs:sequence>
</xs:complexType>
\end{lstlisting}

\section{Exemples d'échanges}
Les échanges suivants ont été obtenus sur la plateforme en fonctionnement (fichiers \code{docs/soap/}).

\begin{lstlisting}[language=XML, caption={Requête \code{verifierCompatibilite} avec en-tête WS-Security}]
<soapenv:Envelope xmlns:soapenv="http://schemas.xmlsoap.org/soap/envelope/"
    xmlns:off="http://stages.universite.com/offres"
    xmlns:wsse="http://docs.oasis-open.org/wss/2004/01/oasis-200401-wss-wssecurity-secext-1.0.xsd">
  <soapenv:Header>
    <wsse:Security soapenv:mustUnderstand="1">
      <wsse:UsernameToken>
        <wsse:Username>admin</wsse:Username>
        <wsse:Password>Admin@2026</wsse:Password>
      </wsse:UsernameToken>
    </wsse:Security>
  </soapenv:Header>
  <soapenv:Body>
    <off:verifierCompatibilite>
      <offreId>1</offreId> <filiere>AMOA</filiere> <niveau>INE2</niveau>
      <competences>BPMN, UML, Java</competences>
    </off:verifierCompatibilite>
  </soapenv:Body>
</soapenv:Envelope>
\end{lstlisting}

\begin{lstlisting}[language=XML, caption={Réponse obtenue}]
<S:Envelope xmlns:S="http://schemas.xmlsoap.org/soap/envelope/">
  <S:Body>
    <ns2:verifierCompatibiliteResponse xmlns:ns2="http://stages.universite.com/offres">
      <compatibilite>
        <offreId>1</offreId>
        <compatible>true</compatible>
        <filiereCompatible>true</filiereCompatible>
        <niveauCompatible>true</niveauCompatible>
        <scoreCompetences>50</scoreCompetences>
        <competenceManquante>SQL</competenceManquante>
        <competenceManquante>Gestion de projet</competenceManquante>
      </compatibilite>
    </ns2:verifierCompatibiliteResponse>
  </S:Body>
</S:Envelope>
\end{lstlisting}

\begin{lstlisting}[language=XML, caption={SOAP Fault de validation (\code{ajouterEtudiant} avec données invalides)}]
<S:Fault>
  <faultcode>S:Server</faultcode>
  <faultstring>Données invalides</faultstring>
  <detail>
    <ns3:ServiceFault xmlns:ns3="http://stages.universite.com/commun">
      <code>DONNEES_INVALIDES</code>
      <message>Données invalides</message>
      <detail>Le champ 'matricule' est invalide (attendu : lettres majuscules suivies de chiffres, ex. E2026001)</detail>
      <detail>Le champ 'email' est invalide (attendu : adresse e-mail)</detail>
      <detail>Le champ 'niveau' doit valoir l'une des valeurs [INE1, INE2, INE3]</detail>
    </ns3:ServiceFault>
  </detail>
</S:Fault>
\end{lstlisting}

\section{Implémentation}
\begin{lstlisting}[language=Java, caption={Contrat (extrait) et implémentation de \code{verifierEtudiant}}]
@WebService(name = "EtudiantService", targetNamespace = Namespaces.ETUDIANTS)
public interface EtudiantService {
    @WebMethod(operationName = "verifierEtudiant")
    @WebResult(name = "verification")
    VerificationEtudiant verifierEtudiant(@WebParam(name = "matricule") String matricule)
            throws ServiceException;
    // ...
}

@Override
public VerificationEtudiant verifierEtudiant(String matricule) throws ServiceException {
    SoapSecurity.utilisateur(context);                 // appelant authentifié (WS-Security)
    if (matricule == null || matricule.isBlank()) {
        throw ServiceException.invalide(List.of("Le matricule est obligatoire"));
    }
    return dao.trouver(matricule)
            .map(e -> {
                boolean eligible = Etudiant.ACTIF.equals(e.getStatut());
                String msg = eligible ? "Étudiant actif, éligible aux candidatures"
                        : "Étudiant non éligible (statut " + e.getStatut() + ")";
                return new VerificationEtudiant(matricule, true, e.getStatut(), eligible, msg);
            })
            .orElseGet(() -> new VerificationEtudiant(matricule, false, null, false,
                    "Aucun étudiant ne correspond au matricule " + matricule));
}
\end{lstlisting}
